\section{Related Work}
\label{sec:related}

Four literatures bear most directly on boundary regret: residual human expertise, deferral and selective prediction, realized human--AI collaboration, and the decision-theoretic value of information.

\paragraph{Residual Human Expertise Audits.}
A line of work asks whether human judgments carry predictive information an algorithm has missed.
\citet{AlurEtAl2023} test for residual expertise using conditional covariance, log-loss improvement, and conditional-independence checks; \citet{AlurRaghavanShah2024} give an algorithmic-indistinguishability account of the same question.
\citet{KleinbergEtAl2018} formalize the selective-labels problem: when human labels are observed only on reviewed cases, estimating how outcomes depend on both the instance and the human judgment requires accounting for non-random review.
\citet{GuoEtAl2025} push the question to the feature level via ILIV-SHAP, attributing each input feature's contribution to the agent-complementary VoI of an AI prediction.
What these papers ask is whether humans carry residual predictive information, and which features drive it. We ask something narrower and decision-relevant: whether that information actually changes the reward-optimal action. Residual expertise can sit comfortably alongside zero residual decision value, and a feature-level decomposition does not tell you whether any of its predictive contributions cross a reward-induced decision facet.

\paragraph{Deferral, Selective Prediction, and Reject Options.}
Learning-to-defer methods train systems to choose between acting automatically and handing an instance to a human expert \citep{MadraseEtAl2018,MozannarSontag2020,HemmerEtAl2023,MaoMohri2023}.
Selective-prediction and reject-option methods address the related problem of when a classifier should abstain under a specified loss \citep{RamaswamyEtAl2018,FrancPrusa2023}, and cost-sensitive variants fold asymmetric penalties into the deferral objective \citep{AlvesEtAl2024}.
All of these decide \emph{when} to defer or abstain given a loss, an expert channel, and a collaboration protocol. The question we want to answer comes earlier: whether the human signal can move the deployed action across a reward-induced boundary at all, before any such protocol is on the table.

\paragraph{Realized Human--AI Collaboration.}
A separate body of work studies when teams actually achieve complementarity in deployment.
\citet{Vaccaro2024} synthesize evidence that human--AI combinations often fail to outperform the better standalone agent, and \citet{PengGargKleinberg2025} prove a structural counterpart: any deterministic strategy combining calibrated forecasts that does not always defer to one agent will sometimes underperform the least accurate agent.
\citet{SchoefferEtAl2024} show that explanations can change reliance without necessarily helping users identify correct advice, while \citet{GuoEtAl2024reliance} propose a decision-theoretic framework that decomposes observed deployment loss into mis-reliance and signal-perception components.
\citet{McLaughlinSpiess2024} model recommendation design for realized complementarity, and \citet{HemmerEtAl2025} organize the empirical evidence on when complementarity appears.
A related critical strand argues that complementarity also turns on action-set design, evaluation metrics, and interaction design \citep{StraitouriEtAl2025}.
That work is about teams once they are deployed: whether they realize value, what protocols are achievable in principle, how observed loss decomposes after the fact. Our question runs upstream of all of that. Is reward-local decision value in the joint distribution to begin with, before strategy choice, behavioral frictions, interfaces, and incentives have any chance to surface it or hide it?

\paragraph{Value of Information, Decision Geometry, and Decision-Focused Learning.}
Conceptually we are closest to value-of-information theory and decision geometry.
\citet{Blackwell1951} orders information structures by decision value. \citet{RadnerStiglitz1984} and \citet{DeLaraGossner2020} characterize when information value vanishes because posteriors stay inside a single action region, and \citet{KamenicaGentzkow2011} expresses information value through convex value functions and induced decision regions.
Decision-focused learning makes the adjacent algorithmic point that prediction should be trained for downstream decisions rather than generic accuracy \citep{Elmachtoub2022,Wilder2019}, and decision calibration asks when predictions are indistinguishable to downstream best-response rules \citep{Zhao2021}.
In clinical statistics, decision curve analysis \citep{Vickers2006} evaluates predictive models under threshold-encoded reward rather than generic accuracy, plotting net benefit across a range of clinician decision thresholds; that work established the principle that model utility should be assessed under its downstream loss.
All of this is about decisions in the abstract, or about how a single model performs under reward. Our move is to take the same geometry and apply it as a human--AI audit: a reward-local check on whether the human signal has decision value for the action actually being deployed, with the underlying VoI identity verified in Lean.
